% Subsection 4.4.1: Zero-shot CoT 流程

Zero-shot CoT（CoT-ZS）对应 H2/H5/H8 三组推理侧。其核心思想是显式切分「推理」与「编码」两步，各自用专属 system 指令约束输出形态，图~\ref{fig:ch4-cot-pipeline} 给出整体数据流。

\begin{figure}[!htb]
  \centering
  \begin{tikzpicture}[
    node distance=5mm and 6mm,
    box/.style={draw, rounded corners=2pt, align=center, minimum height=8mm, minimum width=26mm, font=\footnotesize},
    lbl/.style={font=\scriptsize\itshape, align=center},
    arr/.style={-{Stealth[length=2mm]}, thick},
  ]
    \node[box] (q) {题目 $q$};
    \node[box, right=of q] (s1) {Stage-1\\ Reasoning LLM};
    \node[box, right=of s1] (r) {\texttt{reasoning}};
    \node[box, right=of r] (s2) {Stage-2\\ Code LLM};
    \node[box, right=of s2] (c) {\texttt{solution}};

    \draw[arr] (q) -- (s1);
    \draw[arr] (s1) -- (r);
    \draw[arr] (r) -- (s2);
    \draw[arr] (q.south) to [bend right=14] (s2.south);
    \draw[arr] (s2) -- (c);

    \node[lbl, below=1mm of s1] {SYSTEM: algorithm teacher};
    \node[lbl, below=1mm of s2] {SYSTEM: Python programmer};
    \node[lbl, below=1mm of r]  {截断 + 长度预算};
  \end{tikzpicture}
  \caption{Zero-shot CoT 两阶段推理数据流}
  \label{fig:ch4-cot-pipeline}
\end{figure}

\paragraph{（1）Stage-1：推理生成}
第一阶段只负责产出自然语言推理，不允许直接给出代码：System 为「\texttt{expert algorithm teacher. Generate step-by-step reasoning.}」；User 采用与 \texttt{LoRA\_cot} 训练相同的 CoT 版题目指令（题面 + 末尾 \texttt{Think step by step, then output the Python code.}），\textbf{不是} 4.2.1 节 Direct 模板中的 \texttt{Only output the Python code, no explanation.}，以避免「要求只输出代码」与「本阶段只生成推理」相矛盾；采样参数 \texttt{MAX\_TOKENS\_REASONING=1024}、\texttt{temperature=0.3}、\texttt{top\_p=0.95}、\texttt{top\_k=50}，与 Direct 一致。输出清洗：若模型提前滑出任意三反引号代码块（\verb|```python|、\verb|```bash| 等任意语言标记或无标记围栏），立刻截断只保留前置自然语言段，避免「推理阶段偷跑代码」污染第二阶段上下文。逐字模板见附录~\ref{appendix:prompts}~A.3。

\paragraph{（2）推理段的长度预算}
第一阶段输出不能无限拼进第二阶段，否则会压缩 Stage-2 的生成预算、并在 CoT-FS 模式下挤出示例内容。为此引入两级截断：
\begin{itemize}
    \item \textbf{字符级}：\texttt{MAX\_REASONING\_CHARS=1200}，按字符截断长推理；
    \item \textbf{token 级}：\texttt{MAX\_REASONING\_PROMPT\_TOKENS=256}，按 tokenizer 度量再做一次硬截断，保证即使中文/符号密集也能控制在 Stage-2 prompt 中的占位；
    \item \textbf{VLLM 上下文}：\texttt{VLLM\_MAX\_MODEL\_LEN=8192}，相比 Base 的 2048 显著放宽，为两阶段上下文与 few-shot 示例留出空间。
\end{itemize}

\paragraph{（3）Stage-2：代码生成}
第二阶段接收 \{题目, 截断后的推理\} 作为 chat 上下文，要求输出可直接执行的代码：System 为「\texttt{expert Python programmer. Based on the reasoning, write the code.}」；Assistant 期望格式固定「推理文字 + 空行 + \verb|```python ... ```|」，与训练阶段 assistant 回复格式严格一致；采样参数 \texttt{MAX\_TOKENS\_CODE=1024}、其余与 Stage-1 相同。代码抽取：\texttt{eval\_lora\_cot\_passk.py} 中取\textbf{最后一个} \verb|```python| 块（若无则取最后一个无语言围栏块，再退回全文），以便在模型输出多段围栏时仍以最终代码块为准；此逻辑与 4.2.2 节 Direct 路径「取首段围栏」不同，两份 \texttt{extract\_code} 实现仅在 \texttt{[0]} / \texttt{[-1]} 上有差异。逐字 Stage-2 综合 User 模板见附录~\ref{appendix:prompts}~A.3。

\paragraph{（4）端到端评测入口}
CoT-ZS 的完整评测由 \texttt{evaluation/eval\_lora\_cot\_passk.py} 在 \texttt{--few-shot-k 0} 下启动；数据加载、子进程隔离执行、pass@k 无偏估计与 4.2.2 节\textbf{口径}一致，但\textbf{代码抽取规则}见 4.4.1（3）末条（取末段围栏）。实现上将「一次生成 $N{\times}10$ 份结果」拆成「先批量生成 reasoning、再批量生成 code」两次 \texttt{llm.generate} 调用，以匹配两阶段流水线。